\section{Uso de ferramentas de inteligência
artificial}\label{uso-de-ferramentas-de-inteliguxeancia-artificial}

OpenAI Codex \cite{openai_codex} foi utilizado para planejar o pipeline, auxiliar na
implementação e depuração do código, organizar e revisar o texto e
conferir a consistência entre tabelas e resultados. O autor definiu os
alvos e as restrições da pesquisa, examinou os artefatos produzidos e é
responsável pelo código, pelas afirmações e pela interpretação final. A
ferramenta deve ser citada na versão formatada do artigo conforme a
orientação da disciplina.

O \texttt{registro\ de\ uso} resume as tarefas, mas não preserva todos
os prompts e respostas integrais. O \texttt{suplemento\ de\ IA} contém
cinco prompts verbatim disponíveis no histórico, incluindo o pedido
desta revisão, e explicita a incompletude do registro. Não foram
inventados prompts de sessões de implementação ausentes. Erros
relevantes de assistência incluíram risco de transportar conclusões do
pipeline binário para a regressão, confundir associação com causalidade
e tratar a avaliação externa como seleção. As afirmações quantitativas
deste relatório foram reconciliadas com os CSVs e manifestos mencionados
no Apêndice A; a bibliografia ainda requer conferência editorial final
antes da submissão.
